\section{Source baseline and mathematical contributions}
\label{intro:scope}

This article continues the polylogarithm manuscript and incoming
research reports in the ProveIt repository~\cite{ProveIt} at the fixed
revision
\begin{center}\snapshot.\end{center}
The source comparison uses the manuscript, the relevant archived
continuations, and all eleven incoming ZIP archives present at this
revision. The incoming directory is a research staging area: a claim
being present there is not evidence that it is proved or integrated
into a canonical manuscript chapter. We therefore trace each advance
to a named question and preserve the qualifications in its source.

The immediate starting points are the cubic harmonic
bridge~\cite{IncomingCubic}, the centered parity and rational-resolvent
theorems~\cite{IncomingParity}, and the higher-reflection
generators~\cite{IncomingReflection}. The strict ordered-Hurwitz
conventions come from the all-depth germ report~\cite{IncomingOrdered}.
The exact-relation obstruction and the incomplete larger relation
search~\cite{IncomingExact,IncomingIndependent} delimit what can
presently be claimed about the Gaussian polylogarithm conjectures.

\subsection{Three resolved research directions}

The main advances are the following. They are new relative to the
inspected source corpus; no exhaustive claim of priority across the
literature is intended.

\begin{enumerate}
\item \textbf{The explicitly requested quartic harmonic reduction.}
Question 6 of the cubic harmonic report asks for the finite reduction
at $k=4$, with the higher-depth data identified. We prove
Theorem~\ref{q4:main}, fixing the entire remainder in the quartic
Mittag--Leffler expansion and all integration constants. The answer
uses the linear Laurent coefficient $\theta$ of $Z_3(1+u,1,1)$,
the inherited coefficient $\eta$ of $Z_2(1+u,1)$, and ordinary
zeta and Stieltjes constants. Equation~\eqref{q4:single-coordinate}
eliminates $\eta$ by combining the fourth and third moments.
An absolutely convergent harmonic series, an elementary Mellin
kernel for $\theta$, a normalized primitive, and every argument
derivative accompany the evaluation.

\item \textbf{Every moment of every pair of even harmonic tails.}
Question R1 of the harmonic-parity report asks first for the mixed
pair of orders two and four, and then for a controlled algebra at
general orders. Theorems~\ref{tail:all-moments} and
\ref{tail:generator} solve the pure two-tail family for all even
$p,q\ge2$ and all moment orders. The finite formula retains no
unreduced double-zeta coordinate: all double zeta values have odd
weight and reduce by classical identities. The moment order does
not enlarge the specified arithmetic span. For $(p,q)=(2,4)$,
every moment is a rational combination of the first four. The
additional factors of $H_n$ and odd-order harmonic numbers mentioned
in R1 remain outside this theorem.

\item \textbf{Complex powers and full moving-endpoint completion.}
Question R4 of the same report asks for a branch, a joint
meromorphic generator, and the full resonant amplitudes for
$(\pi\cot\pi x-z)^{-\lambda}$. Theorems
\ref{power:eulerthm}, \ref{power:hurwitzthm}, and
\ref{power:completion} provide them. The result applies to every
nonreal $z$, every Stieltjes index and argument derivative, and
every integer power with arbitrary logarithmic powers. Away from
integer powers the uncompleted local formula applies directly.
The correction at negative integers is finite even when its
denominator is removable; omitting it changes the answer.
Section~\ref{moment:sec} adds a Gamma--Gauss generator for all
ordinary polynomial moments and a generalized-harmonic formula
at the imaginary base points.
\end{enumerate}

The status changes are deliberately narrower than the broadest
questions in the sources. The quartic calculation does not settle
every power $k\ge4$. The two-tail theorem does not treat four tails
or extra harmonic factors. A functional transform for arbitrary
complex powers need not reduce every one of its parameter values to
a finite list of familiar constants. Each theorem states exactly
what is finite, what is continued, and what arithmetic coordinate
remains.

\subsection{Representative identities}

Two compact consequences illustrate the range of the results.
With $x_n=n+1/2$, one obtains the ordinary absolutely convergent sum
\begin{equation}\label{intro:rational-example}
\begin{aligned}
\sum_{n=0}^{\infty}\bigg\{&
\frac{1344x_n^6+2000x_n^4+412x_n^2+7}{6}
                  \psi'(n+1)\psi'''(n+1)\\
&\hspace{17mm}-448x_n^2-\frac{1216}{3}\bigg\}=-88.
\end{aligned}
\end{equation}
The polynomial subtraction is part of the summand; it is not a
separate divergent series. The proof is exact elimination among
four independently specified centered moments, as detailed in
\eqref{tail:24-eliminations}--\eqref{tail:rational-polygamma}.

At a negative half power of the resolvent, an ordinary real integral
has the elementary zeta/logarithm evaluation
\begin{equation}\label{intro:half-example}
\int_0^1\left(x^2-x+\frac16\right)
       \frac{\cos(\pi x/2)}{\sqrt{\sin(\pi x)}}\,\dd x
       =\frac1{12}-\frac{\log^22}{\pi^2}.
\end{equation}
Its entire sequence of Bernoulli analogues follows from one Gamma
quotient, rather than separate evaluations of successive moments.

These ordinary identities complement the finite-part identity for
$\psi^4$. A finite part is appropriate in that case because the
unsubtracted integral diverges at zero; it is not required for
\eqref{intro:rational-example} or \eqref{intro:half-example}.

\subsection{Classical inputs and proof architecture}

The classical ingredients are digamma recurrence and reflection,
Hurwitz Fourier expansions, Mellin transforms, Bernoulli
polynomials, Euler sums, and beta/hypergeometric integral formulas
\cite{DLMF,BBG,EspinosaMoll}. Harmonic-zeta and Stieltjes
coordinates have their own established literature; Coppo's
harmonic-zeta identities provide a recent point of comparison
\cite{Coppo}. We use these sources for specified formulas, not to
infer any unproved constant reduction.

Three proof mechanisms perform different jobs. A two-point contour
fixes the entire function that remains after matching the digamma
principal parts. Exact finite summation by parts identifies the
constant at infinity in a harmonic-tail moment without rearranging
a divergent double sum. Finally, local monomial subtraction in two
complex parameters distinguishes a meromorphic spectral value
from the coordinate finite part of its specialized integrand.
The last distinction is essential for differentiation and
antidifferentiation at resonances.

The proofs do not depend on a numerical relation search. The
accompanying scripts check finite algebra independently and compare
analytically different numerical representations. Explicit
truncation estimates support these comparisons, but the recorded
floating-point runs do not constitute interval certificates or
formal proof-assistant proofs.

